% English translation of checked_pages/esq_119_043.tex; source: cote 119, site 43, PDF page 44, typescript page 41.
% AI-translated by GPT-6 Astra (OpenAI, ChatGPT Pro mode).
\scanpage{119/43}
\mspage{41}

questions, an axiom of stability under analytic continuation
(in $\mathbb{C}^{n}$) appears necessary. The most draconian axiom that
I have been led to introduce, also a stability axiom,
concerns the integration of Pfaff systems, ensuring that certain
Lie groups, or even all of them, are $\unsure{\sourceglyph{119A}}$-manifolds. In all
this, \uline{I have taken care not to assume that the} $\unsure{\sourceglyph{119A}}_{n}$ \uline{are}
\uline{$\mathbb{R}$}-\uline{algebras}, so a constant function on an $\unsure{\sourceglyph{119A}}$-manifold is
``admissible'' only when its value belongs to a certain subfield
$K$ of $\mathbb{R}$ (this is, if one wishes, $\unsure{\sourceglyph{119A}}_{0}$). This subfield may very well
be $\mathbb{Q}$, or its algebraic closure $\bar{\mathbb{Q}}_{r}$ in $\mathbb{R}$, or any other
subextension of $\mathbb{R}/\mathbb{Q}$, preferably even of finite,
or at least countable, transcendence degree over $\mathbb{Q}$. This makes it possible,
for example, as earlier for tame spaces, to associate
with every point $x$ of a manifold (of type $\unsure{\sourceglyph{119A}}$) a residue
field $k(x)$, which is a subextension of $\mathbb{R}/K$. A fact that
seems important to me here is that, even in its strongest form,
the system of axioms does \uline{not} imply that one must have $K = \mathbb{R}$. More
precisely, since \uline{all} the axioms are stability
axioms, it follows that for a given set $S$ of germs
of real analytic functions at the origin (in various spaces\frapparatus{V-FR-043-001}
$\mathbb{R}^{n}$), there exists a smallest theory $\unsure{\sourceglyph{119A}}$ for which these germs
are admissible, and that this theory is ``countable'', i.e. the $\unsure{\sourceglyph{119A}}_{n}$
are countable, as soon as $S$ is countable. A fortiori, $K$ is then
countable, i.e. of countable transcendence degree over $\mathrm{Q}$.

The idea here is to introduce, by means of this
axiomatic approach, a notion of an ``\uline{elementary}'' (real analytic) function,
or rather, a whole hierarchy of such notions. For a function
of 0 variables, i.e. a constant, this notion gives that
of an ``elementary constant'', including in particular (in the case of
the strongest system of axioms) constants such as $\pi$, $e$, and
a multitude of others, by taking values of admissible functions
(such as the exponential, the logarithm, etc.) for systems
of ``admissible'' values of the argument. One feels that the relationship
between the system $\unsure{\sourceglyph{119A}} = (\unsure{\sourceglyph{119A}}_{n})_{n\in\underline{N}}$ and the field of rationality $K$

\margin{43}
